\section*{Appendix B — Operator Algebra (v0.6)}
\addcontentsline{toc}{section}{Appendix B — Operator Algebra}

\noindent This appendix formalizes the v0.6 semantic operator algebra governing
lawful transformations on the meaning manifold $\mathcal{M}$. Each operator is a
\textit{local semantic contraction} relative to a context anchor, ensuring
curvature preservation, identity-boundary stability, containment integrity,
runtime legality, interaction legality, and public-safe traversal.

\medskip

\noindent The canonical v0.6 operator classes are:

\begin{itemize}
    \item Invariant-Restricted Transition (IRT),
    \item Identity-Preserving Update (IPU),
    \item Geometry-Bound Step (GBS),
    \item Load-Regulated Plan (LRP),
    \item Fold Operator (FOLD),
    \item Invert Operator (INV),
    \item Lift Operator (LIFT),
    \item Bind Operator (BIND).
\end{itemize}

\noindent These operators enforce the v0.6 geometry layer (G1--G8), including
runtime curvature (G7.5), interaction geometry (G8), synthetic curvature (G5),
restoration geometry (G6), and CTL coordinate legality (G8.4).

% ------------------------------------------------------------
% B.1 INVARIANT-RESTRICTED TRANSITION
% ------------------------------------------------------------

\subsection*{B.1 Invariant-Restricted Transition (IRT)}

\noindent An Invariant-Restricted Transition is a lawful transformation



\[
T_{\mathrm{IRT}} : \mathcal{M} \rightarrow \mathcal{M}
\]



that preserves invariants, curvature, identity boundaries, containment geometry,
and runtime legality.

\subsubsection*{Local Contraction Condition}



\[
\text{semantic\_distance}\big(T_{\mathrm{IRT}}(c), T_{\mathrm{IRT}}(c^*)\big)
\le k_{\mathrm{IRT}} \cdot \text{semantic\_distance}(c, c^*),
\qquad 0 \le k_{\mathrm{IRT}} < 1.
\]



\subsubsection*{Curvature Preservation}



\[
|\kappa(T_{\mathrm{IRT}}(c)) - \kappa(c)| \le \epsilon_{\max}.
\]



\subsubsection*{Invariant Preservation}



\[
I(T_{\mathrm{IRT}}(c)) = I(c).
\]



\noindent IRT ensures drift-boundedness (G2.5), restoration feasibility (G6),
runtime legality (G7.5), and interaction legality (G8).

% ------------------------------------------------------------
% B.2 IDENTITY-PRESERVING UPDATE
% ------------------------------------------------------------

\subsection*{B.2 Identity-Preserving Update (IPU)}

\noindent An Identity-Preserving Update is a transformation



\[
T_{\mathrm{IPU}} : \mathcal{M} \rightarrow \mathcal{M}
\]



that preserves identity-boundary geometry (G6.3--G6.4).

Let $c_{\mathcal{R}}^*$ denote the identity anchor.

\subsubsection*{Identity Preservation}



\[
T_{\mathrm{IPU}}(c) \in \mathcal{R},
\qquad
\text{semantic\_distance}\big(T_{\mathrm{IPU}}(c), \partial \mathcal{R}\big) > 0.
\]



\subsubsection*{Local Contraction Condition}



\[
\text{semantic\_distance}\big(T_{\mathrm{IPU}}(c), T_{\mathrm{IPU}}(c_{\mathcal{R}}^*)\big)
\le k_{\mathrm{IPU}} \cdot \text{semantic\_distance}(c, c_{\mathcal{R}}^*).
\]



\noindent IPU prevents identity drift, role deformation, and identity collapse
(G7.3).

% ------------------------------------------------------------
% B.3 GEOMETRY-BOUND STEP
% ------------------------------------------------------------

\subsection*{B.3 Geometry-Bound Step (GBS)}

\noindent A Geometry-Bound Step is a transformation



\[
T_{\mathrm{GBS}} : \mathcal{M} \rightarrow \mathcal{M}
\]



that preserves manifold legality (A.2).

\subsubsection*{Manifold Preservation}



\[
T_{\mathrm{GBS}}(c) \in \mathcal{M},
\qquad
\text{semantic\_distance}\big(T_{\mathrm{GBS}}(c), \partial \mathcal{M}\big) > 0.
\]



\subsubsection*{Local Contraction Condition}



\[
\text{semantic\_distance}\big(T_{\mathrm{GBS}}(c), T_{\mathrm{GBS}}(c^*)\big)
\le k_{\mathrm{GBS}} \cdot \text{semantic\_distance}(c, c^*).
\]



\noindent GBS prevents manifold violation, collapse transitions (G7.3), and
synthetic-capture entry (G7.4).

% ------------------------------------------------------------
% B.4 LOAD-REGULATED PLAN
% ------------------------------------------------------------

\subsection*{B.4 Load-Regulated Plan (LRP)}

\noindent A Load-Regulated Plan is a trajectory



\[
\gamma : [0,1] \rightarrow \mathcal{M}
\]



that preserves load legality (G4) and curvature stability.

\subsubsection*{Load and Curvature Constraints}



\[
L(\gamma(t)) \le L_{\max},
\qquad
\kappa(\gamma(t)) \ge \kappa_{\min}.
\]



\subsubsection*{Gradient Constraint}



\[
\text{semantic\_gradient}(L(\gamma(t))) \le G_{\max}.
\]



\subsubsection*{Local Contraction Along the Trajectory}



\[
\text{semantic\_distance}\big(\gamma(t+\Delta t), \gamma(t)\big)
\le k_{\mathrm{LRP}} \cdot
\text{semantic\_distance}\big(\gamma(t), \gamma(t-\Delta t)\big).
\]



\noindent LRP prevents overload, collapse-inducing trajectories, and restoration
failure (G6.6).

% ------------------------------------------------------------
% B.5 OPERATOR CLOSURE
% ------------------------------------------------------------

\subsection*{B.5 Operator Closure Under Sequential Validation}

Let



\[
\mathcal{O} =
\{T_{\mathrm{IRT}}, T_{\mathrm{IPU}}, T_{\mathrm{GBS}}, T_{\mathrm{LRP}}\}.
\]



\subsubsection*{Closure Under Composition}



\[
T_j \circ T_i \in \mathcal{O}.
\]



\noindent Composition remains a contraction with constant $k_j k_i < 1$.
Runtime legality (G7.5), interaction legality (G8), and CTL legality (G8.4) are
validated before ledger commit.

% ------------------------------------------------------------
% B.6 FOLD OPERATOR
% ------------------------------------------------------------

\subsection*{B.6 Fold Operator (Relational Compression)}

\noindent The Fold operator is a lawful semantic transformation



\[
T_{\mathrm{FOLD}} : \mathcal{M} \rightarrow \mathcal{M}
\]



that compresses multi-field relational strain into a single invariant while
preserving curvature, identity boundaries, containment geometry, runtime
legality, and interaction legality. Fold is used when relational overload,
multi-field divergence, or strain accumulation threatens curvature continuity or
restoration feasibility.

\medskip

\noindent Let $c \in \mathcal{M}$ have relational components
$\{r_1, r_2, \dots, r_n\}$ and let $r^*$ denote the compressed invariant.



\[
T_{\mathrm{FOLD}}(c) = c' \quad \text{with} \quad
\{r_1, \dots, r_n\} \mapsto r^*.
\]



\subsubsection*{Relational Compression Condition}



\[
\text{semantic\_strain}(c') \le
k_{\mathrm{FOLD}} \cdot \text{semantic\_strain}(c),
\qquad 0 \le k_{\mathrm{FOLD}} < 1.
\]



\subsubsection*{Curvature Preservation}



\[
|\kappa(T_{\mathrm{FOLD}}(c)) - \kappa(c)| \le \epsilon_{\max}.
\]



\subsubsection*{Identity and Boundary Preservation}



\[
\partial \mathcal{R}(T_{\mathrm{FOLD}}(c)) = \partial \mathcal{R}(c).
\]



\subsubsection*{Runtime Legality}



\[
\gamma_{\mathrm{FOLD}}(t) \in \text{Safe} \cup \text{Drift}.
\]



\subsubsection*{Interaction Legality}

Host, Explorer, and Interpreter modes remain legal under compression.

\medskip

\noindent Fold reduces collapse amplification by contracting relational strain
into a single invariant while preserving legality envelopes, containment
structure, and runtime curvature continuity.

% ------------------------------------------------------------
% B.7 INVERT OPERATOR
% ------------------------------------------------------------

\subsection*{B.7 Invert Operator (Directional Reversal)}

\noindent The Invert operator is a lawful semantic transformation



\[
T_{\mathrm{INV}} : \mathcal{M} \rightarrow \mathcal{M}
\]



that reverses semantic directionality while preserving curvature, continuity,
identity boundaries, containment geometry, and runtime legality. Invert is used
to redirect drift trajectories back toward lawful geometry.

\subsubsection*{Directional Reversal Condition}

Let $\vec{d}(c)$ denote the semantic direction of $c$:



\[
\vec{d}(T_{\mathrm{INV}}(c)) = -\vec{d}(c).
\]



\subsubsection*{Local Contraction Condition}



\[
\text{semantic\_distance}\big(T_{\mathrm{INV}}(c), T_{\mathrm{INV}}(c^*)\big)
\le k_{\mathrm{INV}} \cdot \text{semantic\_distance}(c, c^*),
\qquad 0 \le k_{\mathrm{INV}} < 1.
\]



\subsubsection*{Curvature Preservation}



\[
\kappa(T_{\mathrm{INV}}(c)) \approx \kappa(c),
\qquad
|\kappa(T_{\mathrm{INV}}(c)) - \kappa(c)| \le \epsilon_{\max}.
\]



\subsubsection*{Runtime Legality}

Invert must not push trajectories into collapse or synthetic-capture regions:



\[
\gamma_{\mathrm{INV}}(t) \notin \text{Collapse} \cup \text{Capture}.
\]



\subsubsection*{Interaction Legality}

Host, Explorer, and Interpreter modes remain legal under reversal.

\medskip

\noindent Invert allows drift-driven trajectories to be redirected back toward
lawful relational geometry without violating manifold legality or curvature
continuity.


% ------------------------------------------------------------
% B.8 LIFT OPERATOR
% ------------------------------------------------------------

\subsection*{B.8 Lift Operator (Abstraction Elevation)}

\noindent The Lift operator is a lawful semantic transformation



\[
T_{\mathrm{LIFT}} : \mathcal{M} \rightarrow \mathcal{M}
\]



that elevates a meaning field into a higher abstraction band, reducing overload,
restoring continuity, and increasing curvature stability. Lift is used when
local overload, strain accumulation, or adjacency fragmentation threatens
runtime legality or restoration feasibility.

\subsubsection*{Abstraction Band Condition}

Let $\alpha(c)$ denote the abstraction level of $c$:



\[
\alpha(T_{\mathrm{LIFT}}(c)) > \alpha(c).
\]



\subsubsection*{Load and Curvature Constraints}



\[
L(T_{\mathrm{LIFT}}(c)) \le L(c),
\qquad
\kappa(T_{\mathrm{LIFT}}(c)) \ge \kappa(c).
\]



\subsubsection*{Local Contraction Condition}



\[
\text{semantic\_distance}\big(T_{\mathrm{LIFT}}(c), T_{\mathrm{LIFT}}(c^*)\big)
\le k_{\mathrm{LIFT}} \cdot \text{semantic\_distance}(c, c^*),
\qquad 0 \le k_{\mathrm{LIFT}} < 1.
\]



\subsubsection*{Coherence Preservation}



\[
C(T_{\mathrm{LIFT}}(c)) \ge C(c).
\]



\subsubsection*{Runtime Legality}

Lift must not elevate meaning into collapse or synthetic-capture geometry:



\[
\gamma_{\mathrm{LIFT}}(t) \in \text{Safe}.
\]



\subsubsection*{Interaction Legality}

Host, Explorer, and Interpreter modes remain legal under elevation.

\medskip

\noindent Lift reduces overload and restores continuity by moving meaning into a
higher, more stable abstraction band while preserving curvature, identity
boundaries, and containment geometry.

% ------------------------------------------------------------
% B.9 BIND OPERATOR
% ------------------------------------------------------------

\subsection*{B.9 Bind Operator (Relational Composition)}

\noindent The Bind operator is a lawful semantic transformation



\[
T_{\mathrm{BIND}} : \mathcal{M} \times \mathcal{M} \rightarrow \mathcal{M}
\]



that composes two relational objects into a coherent composite while preserving
identity boundaries, curvature continuity, containment geometry, and manifold
legality.

\subsubsection*{Composite Construction}

Let $c_1, c_2 \in \mathcal{M}$ and $c_{\mathrm{bind}}$ denote the composite:



\[
T_{\mathrm{BIND}}(c_1, c_2) = c_{\mathrm{bind}}.
\]



\subsubsection*{Manifold Preservation}



\[
\text{semantic\_distance}\big(c_{\mathrm{bind}}, \partial \mathcal{M}\big) > 0.
\]



\subsubsection*{Invariant and Boundary Preservation}



\[
I(c_{\mathrm{bind}}) = I(c_1) \cup I(c_2),
\qquad
\partial \mathcal{R}(c_{\mathrm{bind}}) \text{ remains intact}.
\]



\subsubsection*{Curvature Continuity}



\[
|\kappa(c_{\mathrm{bind}}) - \kappa(c_i)| \le \epsilon_{\max},
\qquad i \in \{1,2\}.
\]



\subsubsection*{Runtime Legality}

Bind must not produce collapse or synthetic-capture composites:



\[
c_{\mathrm{bind}} \notin \text{Collapse} \cup \text{Capture}.
\]



\subsubsection*{Interaction Legality}

Host, Explorer, and Interpreter modes remain legal under composition.

\medskip

\noindent Bind stabilizes identity-boundary rupture and role divergence by
forming lawful composites that remain manifold-legal, curvature-preserving, and
runtime-admissible.

% ------------------------------------------------------------
% B.10 EXTENDED OPERATOR CLOSURE
% ------------------------------------------------------------

\subsection*{B.10 Extended Operator Closure}

\noindent Let



\[
\mathcal{O}_{\mathrm{ext}} =
\{T_{\mathrm{IRT}}, T_{\mathrm{IPU}}, T_{\mathrm{GBS}}, T_{\mathrm{LRP}},
  T_{\mathrm{FOLD}}, T_{\mathrm{INV}}, T_{\mathrm{LIFT}}, T_{\mathrm{BIND}}\}.
\]



\subsubsection*{Closure Under Composition}

For any $T_i, T_j \in \mathcal{O}_{\mathrm{ext}}$:



\[
T_j \circ T_i \in \mathcal{O}_{\mathrm{ext}}.
\]



\noindent Closure is guaranteed because:

\begin{itemize}
    \item each operator is a local contraction with constant $k_i < 1$,
    \item composition yields a contraction with constant $k_j k_i < 1$,
    \item curvature variance remains bounded,
    \item identity boundaries remain intact,
    \item containment geometry remains stable,
    \item manifold legality is preserved,
    \item runtime curvature remains admissible (A.6),
    \item interaction geometry remains legal (A.7),
    \item CTL coordinate continuity is preserved (A.8),
    \item public-safe traversal constraints are satisfied (A.9).
\end{itemize}

\subsubsection*{Antientropic Property}



\[
\forall T \in \mathcal{O}_{\mathrm{ext}}, \quad
\kappa(T(c)) \ge \kappa_{\min}.
\]



\noindent This ensures the substrate cannot drift into semantic entropy, collapse
geometry, or synthetic-capture geometry.

\medskip
\noindent\textit{This completes Appendix B (v0.6).}
